% Generated by Listas/make_statements.py from pratica-lecture-08.md, item 15. Do not edit.
\textbf{Source.} Tadelis (2013), \emph{Game Theory: An Introduction}, Exercise 16.6 ``Publish or Perish'', pp.~337--338 {[}PDF pp.~356--357{]}. Reproduced verbatim (with the book's footnote 6). Answer: none in the book.



\begin{quote}

\textbf{Publish or Perish:}⁶ Imagine that any newly minted Ph.D.~who starts a tenure-track assistant professor job (player 1) is one of two types: high-ability (\(\theta_H\)) or low-ability (\(\theta_L\)), where \(\theta_H > \theta_L > 0\). The assistant professor knows his type, but the department that hires him (player 2) knows only that he has high ability with probability \(p < \frac{1}{2}\). The assistant professor first chooses how hard to work, which is effectively how many papers to publish in period 1 (the pre-tenure period). After observing how many papers the assistant professor published, the department decides whether to grant him tenure (\(T\)) or not to do so (\(N\)). If the department chooses to grant tenure then the assistant professor's payoff is \(v_1(q, T \mid \theta) = V - \frac{q}{\theta}\), where \(V\) is the value of being tenured. The department's payoff is 1 if it tenures a high-ability type and \(-1\) if it tenures a low-ability type. If the department denies tenure, it gets a payoff of 0 and the assistant professor's payoff is \(v_1(q, N \mid \theta) = -\frac{q}{\theta}\). Denote by \(\mu(q)\) the department's belief that the professor is a high-ability type given that he published \(q\) papers.



\begin{enumerate}

\def\labelenumi{\alph{enumi}.}

\item

  Write down a ``rough'' extensive form of this game.

\item

  If there is a pooling perfect Bayesian equilibrium, will the assistant professor be tenured? Does he write any papers? What then is the unique outcome of the pooling perfect Bayesian equilibrium?

\item

  Characterize the perfect Bayesian separating equilibria using the parameters of the game.

\item

  If there is more than one separating perfect Bayesian equilibrium, what is the smallest number of papers that an assistant professor needs to publish in order to get tenure?

\end{enumerate}



⁶This exercise is based on one written by Rod Garrett and John Wooders.

\end{quote}

